\section{Python}

\begin{frame}
    \frametitle{本节内容和前置知识}

    本节内容：
    \begin{itemize}
        \item Python
        \item Numpy
        \item PyTorch
        \item Git和GitHub
    \end{itemize}

    默认同学们已经：
    \begin{itemize}
        \item 环境里已经安装了Python，且已经安装了PyTorch。
        \item 会C（基本）。
    \end{itemize}

\end{frame}

\begin{frame}[fragile]
    \frametitle{Python基本使用}

    在命令行中输入 \mintinline{bash}|python| ，就进入了Python的交互式环境，可以直接输入Python代码进行测试和调试。

    输入 \mintinline{bash}|exit()| 或按Ctrl+D可以退出Python交互式环境。

    如欲运行某代码 \verb|code.py| 文件，可以在命令行中输入 \mintinline{bash}|python code.py| 来执行该文件中的代码。

    若你安装了Code，你可以在Code中打开code.py文件，然后按F5或点击运行按钮来运行该文件中的代码。

    你可以在代码的行号处打断点，然后按F5或点击调试按钮来调试代码，这样就可以逐行执行代码，查看变量的值，分析代码的执行流程了。

\end{frame}

\begin{frame}[fragile]
    \frametitle{变量声明}

    Python是一种动态类型语言。虽然依然需要遵循先声明、再使用的原则，但不需要声明变量的类型，变量的类型会在运行时自动推断。

    行尾没有分号，代码块通过缩进来表示，通常使用4个空格的缩进。\textbf{Python中的缩进非常重要，必须使用一致的缩进方式，否则会导致语法错误。}
    \begin{minted}{python}
x = 10 # 实际上这就是个声明，声明x并赋值。
print(x)  # 输出: 10
x = "Hello, World!"
print(x)  # 输出: Hello, World!
    \end{minted}

\end{frame}

\begin{frame}
    \frametitle{基本类型}

    虽然Python是动态类型语言，但它是强类型的。“动态类型”指的是某一变量的类型在运行时确定，换言之，不同时间同一变量的类型可以不同；但一个确定的时间点上，某一变量的类型是确定的，不会同时具有多种类型。

    Python的内置类型包括：int、float、str、bool、list、tuple、dict、set、None等。强制类型转换也在Python中是支持的，例如 \mintinline{python}{int("123")} 会将字符串 "123" 转换为整数 123，无需调用 \mintinline{cpp}|std::stoi| 这样的函数。

    Python3.8以后允许写类型注释，即 \mintinline{python}|x: int = 10| ，但这\textbf{只是注释，并不会真正限制变量的类型}。

\end{frame}

\begin{frame}[fragile]
    \frametitle{输入输出}

    Python中输入输出的基本函数是 \mintinline{python}{input()} 和 \mintinline{python}{print()} 。

    \begin{minted}{python}
name = input("Your Name: ")  # 这会在命令行中显示提示信息，并等待用户输入，用户输入的内容会被赋值给变量name
print("Hello, ", name)  # 输出: Hello, <用户输入的内容>
print("Hello, ", name, end="")  # 不添加换行符
    \end{minted}
    需要注意的是， \mintinline{python}{input()} 函数会将用户输入的内容\textbf{作为字符串}返回，如果需要将输入转换为其他类型（如整数或浮点数），可以使用相应的类型转换函数，如 \mintinline{python}{int()} 或 \mintinline{python}{float()} 。
    
\end{frame}

\begin{frame}[fragile]
    \frametitle{Python的变量运算}

    所以在C++中的加减乘除（单双目）、与或非、大小等、赋值在Python中也适用，但Python还支持一些额外的运算符，如幂运算符 \mintinline{python}{**} 、取模运算符 \verb|%| 、整除运算符 \mintinline{python}{//} 等。

    需要注意的是，在Python中，5/2==2.5，而不是2。这是因为Python中的除法运算符 \mintinline{python}{/} 总是返回一个浮点数结果。如果你想要像C++一样得到整数结果，可以使用整除运算符 \mintinline{python}{//} ，即5//2==2。

    Python没有自增自减运算符，也没有取地址运算符。这是因为Python没有指针，变量的本质都是引用。

    Python用井号 \mintinline{python}{\#} 来表示注释，直接在代码行前加 \mintinline{python}{\#} 即可。

\end{frame}

\begin{frame}[fragile]
    \frametitle{Python条件分支}

    Python中的条件语句使用 \mintinline{python}{if} 、\mintinline{python}{elif} 和 \mintinline{python}{else} 来实现，功能和使用基本上与C中的条件语句相同，但Python使用冒号和缩进来表示代码块，而不是花括号。

    \begin{minted}{python}
if x == 2: # 必须换行并缩进
    print("x is 2")
elif x > 2: # 必须一致地使用缩进
    print("x is greater than 2")
else:
    print("x is less than 2")
    \end{minted}

\end{frame}

\begin{frame}[fragile]
    \frametitle{Python循环}

    Python中循环有两种： \mintinline{python}{while} 循环（和C中的while类似）和 \mintinline{python}{for} 循环（和C++中的range-for类似）。
    \begin{minted}{python}
while x > 0:    # 下一行必须缩进
    print(x)
    x -= 1  # 整个缩进块内都是循环体
    \end{minted}
\end{frame}

\begin{frame}[fragile]
    \frametitle{Python迭代}

    Python的for循环可以直接迭代任何可迭代对象，如列表、字符串、字典等，非常方便。
    \begin{minted}{python}
for i in range(5):  # range(5)生成一个可迭代对象，包含0到4的整数
    print(i)
for i in [1, 2, 3]:  # 直接迭代一个列表
    print(i)
for char in "Hello":  # 直接迭代一个字符串
    print(char) # 输出: H e l l o，每行一个字符
    \end{minted}
    需要说明的是， \mintinline{python}|range(5)| 实际上是一个生成器对象，指的是0到5的整数序列（不包含5），而不是一个列表。
\end{frame}

\begin{frame}[fragile]
    \frametitle{复合数据类型}
    Python中的复合数据类型包括列表（list）、元组（tuple）、字典（dict）和集合（set）。这些数据类型可以存储多个值，并且支持各种操作，如添加、删除、访问等。

    Python的这些复合数据类型都是可以直接打印的，且打印出来的格式非常友好，方便调试和展示。

    \begin{minted}{python}
l: list[int] = [1, 2, 3]  # 列表，有序可更改
t: tuple[int, int, int] = (1, 2, 3)  # 元组，有序不可更改
d: dict[str, int] = {"a": 1, "b": 2}  # 字典，键值对，无序可更改
s: set[int] = {1, 2, 3}  # 集合，无序不重复可更改
    \end{minted}

\end{frame}

\begin{frame}[fragile]
    \frametitle{复合数据类型的操作}

    操作类似C++的STL。以list为例：
    \begin{minted}{python}
l[0] = 10  # 访问并修改元素
l[-1] = 20  # 负数索引表示从末尾开始访问
# l[114514] = 30  # 访问越界依然会报错
l.append(4)  # 末尾添加元素
l.remove(2)  # 删除元素，如果有重复元素只删除第一个
l.sort()  # 排序
l.reverse()  # 反转
l.clear()  # 清空
l.insert(1, 5)  # 在索引1处插入元素5
    \end{minted}

\end{frame}

\begin{frame}[fragile]
    \frametitle{切片}

    切片可以用来获取列表、元组、字符串等序列的一部分。
    
    切片的语法是 \mintinline{python}{sequence[start:stop:step]} ，其中 \mintinline{python}{start} 是起始索引， \mintinline{python}{stop} 是结束索引（不包含）， \mintinline{python}{step} 是步长（可以为负数）。三个参数都是可选的，但不能全部省略。

    例如 \mintinline{python}{l[1:]} 表示从索引1开始，步长为1，获取列表l的所有元素； \mintinline{python}{l[1:4:2]} 表示从索引1开始，到索引4结束（不包含），步长为2，获取列表l的元素； \mintinline{python}{l[::-1]} 表示从列表l的末尾开始，步长为-1，获取列表l的所有元素并反转。切片也允许负索引（含义是从末尾开始访问），但不允许越界访问；也必须保证起始索引比结束索引更小，否则会得到一个空列表。

    切片会返回一个新的序列，原序列不受影响。

\end{frame}

\begin{frame}
    \frametitle{推导式}

    列表推导式是一种简洁的语法，可以用来创建新的列表。它的基本语法是 \mintinline{python}{[expression for item in iterable if condition]} ，其中 \mintinline{python}{expression} 是生成新列表元素的表达式， \mintinline{python}{item} 是迭代器中的每个元素， \mintinline{python}{iterable} 是一个可迭代对象， \mintinline{python}{condition} 是一个可选的条件表达式。

    实际上，推导式不仅可以用来生成列表，还可以用来生成其他类型的集合，如字典推导式和集合推导式。但因为推导式会对可读性造成一定的影响，所以不建议使用过度复杂的推导式。

    推导式的本质实际上是循环。

\end{frame}

\begin{frame}[fragile]
    \frametitle{字符串和f字符串}

    字符串是Python中的一种基本数据类型，用于表示文本数据。字符串可以使用单引号（'）或双引号（"）来定义，也可以使用三引号（'''或"""）来定义多行字符串。

    f字符串（f-string）是Python 3.6引入的一种字符串格式化方法，允许在字符串中直接嵌入表达式，并在运行时进行求值。f字符串以字母f或F开头，后跟一个字符串字面量，表达式用花括号{}括起来。例如：

\end{frame}

\begin{frame}[fragile]
    \frametitle{f-string 实例}

    \begin{minted}{python}
name = "Alice"
age = 30
greeting = f"Hello, my name is {name} and I am {age} years old."
print(greeting)  # 输出: Hello, my name is Alice and I am 30 years old.
    \end{minted}


\end{frame}

\begin{frame}[fragile]
    \frametitle{函数}

    Python中函数的定义使用 \mintinline{python}{def} 关键字，后跟函数名和参数列表。例如：
    \begin{minted}{python}
def add(a: int, b: int) -> int: # 这里的三个int都是类型注释，实际并不限制参数和返回值的类型。
# def add(a, b):  # 也可以不写类型注释
    return a + b
    \end{minted}
    Python中的函数参数可以有默认值。函数也可以返回多个值，实际上是返回一个元组。

\end{frame}

\begin{frame}[fragile]
    \frametitle{可变参数和关键字参数}

    函数参数也可以是可变参数（*args）或关键字参数（**kwargs）。

    可变参数指的是函数参数的数量不固定，可以接受任意数量的参数。换言之，args是一个元组，但可以在传参时不显式地传入一个元组，而是直接传入多个参数，Python会自动将它们打包成一个元组。

    关键字参数类似，指的是函数参数以键值对的形式传递，可以接受任意数量的键值对参数。换言之，kwargs是一个字典，同样可以不显式地传入一个字典，而是直接传入多个键值对参数。

\end{frame}

\begin{frame}[fragile]
    \frametitle{可变参数和关键词参数 (2)}

\begin{minted}{python}
def my_function(*args, **kwargs):
    print("args:", args)
    print("kwargs:", kwargs)
my_function(1, 2, 3, name="Alice", age=30)
\end{minted}

\end{frame}

\begin{frame}[fragile]
    \frametitle{面向对象}

    可以使用 \mintinline{python}{class} 关键字来定义一个类，然后创建该类的实例（对象）来使用它。例如：
    \begin{minted}{python}
class Person:
    def __init__(self, name: str, age: int):
        self.name = name  # 属性
        self.age = age
    def greet(self): # 方法
        print(f"Hello, my name is {self.name} and I am {self.age} years old.")
    \end{minted}

\end{frame}

\begin{frame}[fragile]
    \frametitle{类的实例化}

    可以通过创建 \mintinline{python}{Person} 类的实例来使用上文中提到的类，例如：
    \begin{minted}{python}
person1 = Person("Alice", 30)
person1.greet()  # 输出: Hello, my name is Alice and I am 30 years old.
    \end{minted}
    故而我们刚刚这种调用 \mintinline{python}|l.append(4)| 的方式，其实就是在调用list类的append方法。

\end{frame}

\begin{frame}[fragile]
    \frametitle{类的常用属性}

    一般的，要实现一个类，我们往往需要实现以下常用属性：
    \begin{itemize}
        \item \mintinline{python}|__init__(self, ...)|：构造函数，用于初始化对象的属性。
        \item \mintinline{python}|__str__(self)|：字符串表示方法，用于定义对象的字符串表示形式，方便打印和调试。
        \item \mintinline{python}|__repr__(self)|：对象表示方法，用于定义对象的详细表示形式，通常在调试时使用。
        \item \mintinline{python}|__eq__(self, other)|：相等性比较方法，用于定义对象之间的相等性比较逻辑。
    \end{itemize}
    当然，你也可以根据需要实现其他的特殊方法，如 \mintinline{python}|__add__(self, other)| 用于定义对象之间的加法操作， \mintinline{python}|__len__(self)| 用于定义对象的长度等。Python的运算符重载实际上都是通过实现这些特殊方法来实现的。

\end{frame}

\begin{frame}[fragile]
    \frametitle{数据类}

    Python 3.7引入了数据类（dataclass）的概念，相当于简单的结构体。Python会自动实现其常用方法。

    \begin{minted}{python}
from dataclasses import dataclass

@dataclass # 装饰器，告诉Python这是一个数据类
class Point:
    x: float
    y: float
    \end{minted}

\end{frame}

\begin{frame}[fragile]
    \frametitle{模块和包}

    \begin{minted}{python}
import math  # 相当于#include <cmath>;
from math import sqrt  # 相当于#include <cmath>和using std::sqrt;
from math import *  # 相当于#include <cmath>和using namespace std;
    \end{minted}
    需要注意的是，第三种方式可能会导致命名冲突，因此不建议使用这种方式导入模块。

    允许为模块或函数指定别名，例如 \mintinline{python}{import numpy as np} ，这样就可以使用 \mintinline{python}{np} 来代替 \mintinline{python}{numpy} ，使代码更简洁。


\end{frame}
